% ============================================================
% f2014_shanghai_arts_q21_billboard.tex
% 2014 上海卷（文科）第 21 题：广告牌 CD 的仰角示意图
% 来源 PDF：普通高考/2014/2014上海文.pdf
% 原页：第 2 页（A3），印面 918，第 21 题题干下方（左栏）
% A、C、B 在同一水平线上，AC=35 米、CB=80 米，D 为广告牌顶端；
% 图中 CD 为铅垂方向（第 (1) 问的情形），α、β 分别为从 A、B
% 看 D 的仰角。示意图按 AC:CB≈35:80 的比例绘制，高度取示意值。
% ============================================================
\begin{tikzpicture}[>=Stealth, line width=0.65pt, font=\small, scale=0.95,
  line cap=round, line join=round]
  \coordinate (A) at (0,0);
  \coordinate (C) at (3.50,0);
  \coordinate (B) at (11.50,0);
  \coordinate (D) at (3.50,3.20);

  % 水平线、广告牌与两条视线
  \draw (A) -- (B);
  \draw (D) -- (C);
  \draw (A) -- (D);
  \draw (B) -- (D);

  % 仰角 α、β 的角标
  \draw (1.05,0) arc (0:42.44:1.05);
  \draw ([shift=(158.13:1.10)]B) arc (158.13:180:1.10);

  % 顶点标注
  \fill (A) circle (0.85pt) node[below left=-1pt] {$A$};
  \fill (C) circle (0.85pt) node[below=-1pt] {$C$};
  \fill (B) circle (0.85pt) node[below right=-1pt] {$B$};
  \fill (D) circle (0.85pt) node[above=-1pt] {$D$};
  \node at (1.60,0.52) {$\alpha$};
  \node at (10.45,0.60) {$\beta$};
\end{tikzpicture}
